\chapter{Conclusion}{~}{~}
\label{CH:CONC}

 In this thesis, we proposed Wavelet-Conditioned Band-Decoupled Spatial Feature Transform (Band-Decoupled SFT), a frequency-domain adaptation framework that injects DT-CWT-derived priors into a frozen pre-trained diffusion U-Net. To effectively address temporal flickering in diffusion-based VSR, we leveraged the near shift-invariance of DT-CWT magnitudes, which remain approximately stable under sub-pixel translations. The complex wavelet coefficients of each low-resolution frame are decomposed into four scales and partitioned into high-frequency and low-frequency bands according to their semantic roles. Therefore, two distinct frequency priors are constructed and routed to semantically appropriate U-Net stages by the designed framework.
 
 To maximize the structural and textural guidance, we designed an asymmetric encoder-decoder injection scheme. High-frequency bands modulate the U-Net decoder for textural detail synthesis, while low-frequency bands modulate the encoder for structural feature extraction. This asymmetric injection demonstrated the significance of aligning frequency components with the hierarchical processing of U-Net while preserving the generative prior. Furthermore, processing magnitude and phase through decoupled branches with trigonometric phase encoding enhanced the stability of frequency conditioning.
 

 
From the experiments, we have proved the proposed asymmetric injection and decoupled magnitude-phase processing can achieve strong temporal consistency. The proposed method achieves a 63\% reduction in tLPIPS (from 41.11 to 15.25) on REDS4 compared to the StableVSR baseline. Moreover, the Frequency Conditioning Encoder requires only 6.22M parameters, approximately 3\% of the ControlNet backbone, while consistent improvements are observed across the Vid4, UDM10, and SPMCS benchmarks when compared to StableVSR \cite{rota2024stablevsr}. 

As future work, a method to extend the frequency-domain conditioning paradigm to a fully flow-free pipeline by replacing the optical flow-based warping in the ControlNet branch is needed. Furthermore, various research on latent-space wavelet analysis and adaptive band partitioning is needed.


    
%%%
